\scanpage{44}
\mspage{-VI-5-}
\hfill n°383

non liés. Comme $a$ est une extrémité de $X$, il existe un indice $i$, par exemple $2$, tel que $a$ ne soit lié à aucun noeud de $X'_i$. On peut alors prendre $X_1=X'_1\cup\{a\}$, $X_2=X'_2$. C.Q.F.D.

On reprend les notations du début.

DEFINITION 2.- \uline{On appelle graphe de} $G$ \uline{un graphe dont les noeuds sont les murs} $H_1,\ldots,H_n$ \uline{de} $C$, \uline{deux éléments} $H_i$ \uline{et} $H_j$ $(i\ne j)$ \uline{étant liés si et seulement si} $(e_i|e_j)\ne0$.

Comme $G$ opère transitivement sur ses chambres (§3, n°1, prop.1), les graphes de $G$ sont tous isomorphes.

Supposons $G$ essentiel. Soit $G=G_1\times\ldots\times G_s$ la décomposition de $G$ en produit de groupes irréductibles. Le graphe $X$ de $G$ est isomorphe à la somme disjointe des graphes $X_i$ des $G_i$. Pour que $G$ soit irréductible, y faut et y suffit que son graphe $X$ soit connexe.

PROPOSITION 1.- \uline{Le graphe} $X$ \uline{de} $G$ \uline{est une forêt}. \uline{Si} $G$ \uline{est irréductible}, \uline{ce graphe est un arbre}.

En effet, en vertu du lemme 5, $X$ n'admet pas de circuit.

3. \uline{Définition des transformations de Coxeter}.

Appelons \uline{chambre ordonnée} relative à $G$ la donnée d'une chambre $C$ et d'une indexation $(H_1,\ldots,H_n)$ de ses murs par $1,2,\ldots,n$.

DEFINITION 3.- \uline{On appelle transformation de Coxeter définie par une chambre ordonnée} $C$, $(H_1,\ldots,H_n)$ \uline{l'élément} $s_{H_1}s_{H_2}\ldots s_{H_n}$ \uline{de} $G$.

PROPOSITION 2.- \uline{Toutes les transformations de Coxeter de} $G$ \uline{sont conjuguées dans} $G$.

Soient $c$ et $c'$ deux transformations de Coxeter. Comme $G$ opère transitivement sur les chambres (§3, n°1, prop.1), on est ramené au cas où elles sont définies par la même chambre sous-jacente $C$, et où elles ne diffèrent donc que par l'ordre des
\mspage{43}
